\section{Probability Theory}

\begin{definition}
	A \textbf{discrete probability space} is a set of elementary events $\Omega = \{\omega_1, \dots, \omega_n \}$.
	Every event $\omega_i$ has a probability $P[\omega_i] \in [0,1]$ such that
	$$\sum_{\omega_i \in \Omega} P[\omega_i] = 1$$
\end{definition}

\begin{remark}
	A finite space where every event has equal probability is a \textbf{Laplace-Space}. For any event $E$:
	$$P[E] = \frac{\vert E \vert}{\vert \Omega \vert}$$
\end{remark}

\begin{definition}
	An \textbf{Event} $E \subseteq \Omega$ has probability $P[E] = \sum_{\omega_i \in E} P[\omega_i]$.
	The complementary event $\bar{E} = \Omega \setminus E$ has probability $P[\bar{E}] = 1 - P[E]$.
\end{definition}

\begin{theorem}
	For two events $A, B \subseteq \Omega$:
	$P[\emptyset] = 0$, $P[\Omega] = 1$, $P[A] \in [0,1]$.
	If $A \subset B$ then $P[A] \le P[B]$.
\end{theorem}

\begin{theorem}
	\textbf{Union Bound}: For events $A_1, \dots, A_n$:
	$$P[\bigcup_{i=1}^n A_i] \leq \sum_{i=1}^n P[A_i]$$
	(Equality iff $A_i$ are pairwise disjoint).
\end{theorem}

\begin{theorem}
	\textbf{Inclusion-Exclusion Principle}: For $n \geq 2$:
	$$P[\bigcup_{i=1}^n A_i] = \sum_{k=1}^n (-1)^{k+1} \sum_{I \subset [n], |I|=k} P[\bigcap_{j \in I} A_j]$$
\end{theorem}

\subsection{Conditional Probability}

\begin{definition}
	For $P[B] > 0$, the \textbf{conditional probability} of $A$ given $B$ is:
	$$P[A \mid B] = \frac{P[A \cap B]}{P[B]}$$
\end{definition}

\begin{remark}
	Always holds: $P[A \mid \Omega] = P[A]$, $P[A \mid A] = 1$, $P[A \mid \bar{A}] = 0$, $P[A \cap B] = P[A \mid B] P[B]$.
\end{remark}

\begin{theorem}
	\textbf{Multiplication-Theorem}: If $P[\bigcap_{i=1}^{n-1} A_i] > 0$:
	$$P[\bigcap_{i=1}^n A_i] = P[A_1] \cdot P[A_2 \mid A_1]$$
	$$\cdot \dots \cdot P[A_n \mid \bigcap_{i=1}^{n-1} A_i]$$
\end{theorem}

\begin{theorem}
	\textbf{Theorem of total probability}: For pairwise disjoint $A_i$ and $B \subseteq \bigcup_{i=1}^n A_i$:
	$$P[B] = \sum_{i=1}^n P[B \mid A_i] P[A_i]$$
\end{theorem}

\begin{theorem}
	\textbf{Bayes' Theorem}: For exactly defined disjoint $A_i$, $B \subseteq \bigcup A_i$, and $P[B] > 0$:
	$$P[A_i \mid B] = \frac{P[B \mid A_i] P[A_i]}{\sum_{j=1}^n P[B \mid A_j] P[A_j]}$$
\end{theorem}

\subsection{Independence}

\begin{definition}
	Two events $A, B$ are \textbf{independent} if:
	$$P[A \cap B] = P[A] P[B] \text{ or } P[A \mid B] = P[A]$$
	Events $A_1, \dots, A_n$ are independent if for all $I \subseteq \{1, \dots, n\}$:
	$$P[\bigcap_{i \in I} A_i] = \prod_{i \in I} P[A_i]$$
\end{definition}

\begin{lemma}
	Let $A_1, \dots, A_n$ be events, $A_i^0 = \bar{A}_i$, $A_i^1 = A_i$. They are independent iff $\forall (s_1,\dots,s_n) \in \{0,1\}^n$:
	$$P[\bigcap_{i=1}^n A_i^{s_i}] = \prod_{i=1}^n P[A_i^{s_i}]$$
\end{lemma}

\begin{lemma}
	If $A, B, C$ are independent, then so are $A \cap B$ and $C$, and $A \cup B$ and $C$.
\end{lemma}

\subsection{Random Variables}

\begin{definition}
	A \textbf{random variable} is a map $X: \Omega \to \mathbb{R}$.
	The \textbf{probability mass function (PMF)} is $f_X(x) = P(X=x)$.
	The \textbf{distribution function (CDF)} is $F_X(x) = P(X \leq x)$.
\end{definition}

\subsubsection{Expectation}

\begin{definition}
	The expectation of $X$ is defined as:
	$$\mathbb{E}[X] \coloneqq \sum_{\alpha \in W_x} \alpha P(X=\alpha)$$
\end{definition}

\begin{theorem}
	For a random variable $X$ with values in $\mathbb{N}_0$:
	$$\mathbb{E}[X] = \sum_{k=1}^\infty P(X \geq k)$$
\end{theorem}

\begin{definition}
	Indicator-variable $X_E$ for $E \subseteq \Omega$:
	$X_E(\omega) = 1$ if $\omega \in E$, else $0$.
	$$\mathbb{E}[X_E] = P(E)$$
\end{definition}

\begin{theorem}
	\textbf{Linearity of Expectation}:
	For random variables $X_1, \dots, X_n$ and $a_1, \dots, a_n, b \in \mathbb{R}$:
	$$\mathbb{E}[\sum_{i=1}^n a_i X_i + b] = \sum_{i=1}^n a_i \mathbb{E}[X_i] + b$$
\end{theorem}

\begin{algorithm}
	#Maximal Stable Set Algorithm
	S[v] = 1/0 with probability p
	For all {u,v} in E:
	if S[u] == 1 and S[v] == 1: S[v] = 0
	Expected size: $\mathbb{E}[\mid S \mid] \geq n*p - m*p^2$
\end{algorithm}

\subsubsection{Variance}

\begin{definition}
	For $X$ with $\mathbb{E}[X] = \mu$, the \textbf{variance} is:
	$$Var[X] \coloneqq \mathbb{E}[(X-\mu)^2]$$
	Standard deviation $\sigma \coloneqq \sqrt{Var[X]}$.
\end{definition}

\begin{theorem}
	$$Var[X] = \mathbb{E}[X^2] - (\mathbb{E}[X])^2$$
\end{theorem}

\begin{lemma}
	For $a, b \in \mathbb{R}$: $Var[aX+b] = a^2 Var[X]$
\end{lemma}

\begin{definition}
	$\mathbb{E}[X^k]$ is the $k$-th \textbf{moment}. $\mathbb{E}[(X-\mathbb{E}[X])^k]$ is the $k$-th \textbf{central moment}.
\end{definition}

\subsubsection{Multiple \& Combined Variables}

\begin{definition}
	RVs $X_1, \dots, X_n$ are \textbf{independent} iff $\forall a_i \in W_{X_i}$:
	$$P[X_1=a_1, \dots, X_n=a_n] = \prod_{i=1}^n P[X_i=a_i]$$
\end{definition}

\begin{theorem}
	If $X, Y$ are independent, $f_{X+Y}(\alpha) = \sum_{\beta} f_X(\beta) f_Y(\alpha-\beta)$.
\end{theorem}

\begin{remark}
	\textbf{Calculation Rules}:
	\begin{itemize}
		\item $\mathbb{E}[X+Y] = \mathbb{E}[X] + \mathbb{E}[Y]$
		\item If $X, Y$ independent: $\mathbb{E}[XY] = \mathbb{E}[X] \mathbb{E}[Y]$
		\item If $X, Y$ independent: $Var[X+Y] = Var[X] + Var[Y]$
	\end{itemize}
	Note $Var[XY] \neq Var[X]Var[Y]$ in general.
\end{remark}

\begin{theorem}
	\textbf{Total Expectation}: For pairwise disjoint events $A_1, \dots, A_n$ covering $\Omega$:
	$$\mathbb{E}[X] = \sum_{i=1}^n \mathbb{E}[X \mid A_i] P[A_i]$$
\end{theorem}

\begin{theorem}
	\textbf{Wald's Equation}: If $N$ and $X$ are independent RVs and $Z \coloneqq \sum_{i=1}^N X_i$ (where $X_i$ are independent copies of $X$):
	$$\mathbb{E}[Z] = \mathbb{E}[N] \mathbb{E}[X]$$
\end{theorem}

\subsection{Distributions}

\begin{definition}
	\textbf{Bernoulli($p$)}: $X \in \{0,1\}$
	$$P[X=1] = p \quad P[X=0] = 1-p$$
	$$\mathbb{E}[X] = p \quad Var[X] = p(1-p)$$
\end{definition}

\begin{definition}
	\textbf{Binomial($n, p$)}: $X \in \{0,\dots,n\}$
	$$P[X=k] = \binom{n}{k} p^k (1-p)^{n-k}$$
	$$\mathbb{E}[X] = np \quad Var[X] = np(1-p)$$
\end{definition}

\begin{definition}
	\textbf{Poisson($\lambda$)}:
	$$P[X=k] = \frac{e^{-\lambda} \lambda^k}{k!}$$
	$$\mathbb{E}[X] = \lambda \quad Var[X] = \lambda$$
\end{definition}

\begin{definition}
	\textbf{Geometric($p$)}:
	$$P[X=k] = p(1-p)^{k-1}$$
	$$\mathbb{E}[X] = \frac{1}{p} \quad Var[X] = \frac{1-p}{p^2}$$
	(Number of throws until first success).
\end{definition}

\begin{definition}
	\textbf{Negative Binomial of order $n$}:
	$$P[X=k] = \binom{k-1}{n-1} p^n (1-p)^{k-n}$$
	(Number of throws until $n$-th success).
\end{definition}

\subsection{Bounding Probability}

\begin{theorem}
	\textbf{Markov's Inequality:} For non-negative RV $X$ and $t \ge 0$:
	$$P[X \geq t] \leq \frac{\mathbb{E}[X]}{t} \quad \text{or} \quad P[X \geq t \mathbb{E}[X]] \leq \frac{1}{t}$$
\end{theorem}

\begin{theorem}
	\textbf{Chebyshev's Inequality:} For any RV $X$ and $t \ge 0$:
	$$P[|X-\mathbb{E}[X]| \geq t] \leq \frac{Var[X]}{t^2}$$
\end{theorem}

\subsubsection{Chernoff Bounds}

Applies to sums of independent Bernoulli RVs $X = \sum X_i$ with $P[X_i=1]=p_i$.
\begin{theorem}
	For $0 \leq \delta \leq 1$:
	\begin{align*}
		P[X \geq (1+\delta)\mathbb{E}[X]] & \leq e^{-\frac{1}{3}\delta^2\mathbb{E}[X]} \\
		P[X \leq (1-\delta)\mathbb{E}[X]] & \leq e^{-\frac{1}{2}\delta^2\mathbb{E}[X]}
	\end{align*}
	For larger deviations ($t \geq 2 e \mathbb{E}[X]$):
	$$P[X \geq t] \leq 2^{-t}$$
\end{theorem}

\subsubsection{Target Shooting (Sampling)}
To estimate $|S|/|U|$, sample $n$ elements uniformly. Let $Y = \frac{1}{n} \sum Y_i$ be the fraction of samples in $S$.
\begin{theorem}
	With given $\delta, \epsilon \ge 0$, if $n \geq \frac{3 |U|}{|S| \epsilon^2} \ln \frac{2}{\delta}$, then
	$$Y \in \left[(1-\epsilon)\frac{|S|}{|U|}, (1+\epsilon)\frac{|S|}{|U|}\right]$$
	with probability $\ge 1-\delta$.
\end{theorem}
